%-----------------------------------LICENSE------------------------------------%
%   This file is part of Mathematics-and-Physics.                              %
%                                                                              %
%   Mathematics-and-Physics is free software: you can redistribute it and/or   %
%   modify it under the terms of the GNU General Public License as             %
%   published by the Free Software Foundation, either version 3 of the         %
%   License, or (at your option) any later version.                            %
%                                                                              %
%   Mathematics-and-Physics is distributed in the hope that it will be useful, %
%   but WITHOUT ANY WARRANTY; without even the implied warranty of             %
%   MERCHANTABILITY or FITNESS FOR A PARTICULAR PURPOSE.  See the              %
%   GNU General Public License for more details.                               %
%                                                                              %
%   You should have received a copy of the GNU General Public License along    %
%   with Mathematics-and-Physics.  If not, see <https://www.gnu.org/licenses/>.%
%------------------------------------------------------------------------------%
\documentclass{article}
\usepackage{amssymb}
\usepackage{amsmath}
\usepackage{xcolor}
\usepackage{float}

\title{18.100P: Midterm}
\date{Spring 2026}

% No indent and no paragraph skip.
\setlength{\parindent}{0em}
\setlength{\parskip}{0em}

\begin{document}
    \maketitle
    \textbf{Foundations / Logic}
    \begin{itemize}
        \item
            Write down the truth table for:
            $P\Rightarrow(Q\Rightarrow\neg{P})$.
        \item
            Is it consistent to accept this expression
            as an axiom of logic? Explain using your truth table.
    \end{itemize}
    \color{blue}
    We can write the truth table as follows.
    \begin{table}[H]
        \centering
        \color{blue}
        \begin{tabular}{cc|cccc}
            $P$
                &
                $Q$
                &
                $\neg{P}$
                &
                $Q\Rightarrow\neg{P}$
                &
                $P\Rightarrow(Q\Rightarrow\neg{P})$\\
            \hline
            0&0&1&1&1\\
            0&1&1&1&1\\
            1&0&0&1&1\\
            1&1&0&0&0
        \end{tabular}
        \caption{Truth Table for $P\Rightarrow(Q\Rightarrow\neg{P})$.}
    \end{table}
    It would be ill-advised to accept this as an axiom of logic. An axiom
    should be something that is intuitively \textbf{true}, but the statement
    $P\Rightarrow(Q\Rightarrow\neg{P})$ may be false. In particular, if $P$ is
    true and $Q$ is true, then $P\Rightarrow(Q\Rightarrow\neg{P})$ is false.
    Axioms are statements that are \textit{always} true.
    \color{black}
    \clearpage
    \textbf{Sequences and Subsequences}
    \begin{itemize}
        \item
            Prove that if $f:\mathbb{R}\rightarrow\mathbb{R}$ is continuous
            at $x_{0}\in\mathbb{R}$, and if $a:\mathbb{N}\rightarrow\mathbb{R}$
            is a convergent sequence with
            $a_{n}\rightarrow{x}_{0}$, then the sequence
            $f(a_{n})$ also converges and $f(a_{n})\rightarrow{f}(x_{0})$.
        \item
            Prove that if $a,b\in\mathbb{R}$ with $a<b$, if
            $f:[a,\,b]\rightarrow\mathbb{R}$ is continuous at each point
            $r\in[a,\,b]$, and if $x:\mathbb{N}\rightarrow[a,\,b]$
            is any sequence, then there is a convergent subsequence of
            $f(x_{n})$.
    \end{itemize}
    \color{blue}
    Let $\varepsilon>0$ be given. To show $f(a_{n})\rightarrow{f}(x_{0})$, we
    must show that there is an $N\in\mathbb{N}$ such that $n\in\mathbb{N}$ and
    $n>N$ implies $|f(a_{n})-f(x_{0})|<\varepsilon$. But $f$ is continuous
    at $x_{0}$, and hence there is a $\delta>0$ such that
    $|x-x_{0}|<\delta$ implies $|f(x)-f(x_{0})|<\varepsilon$. But by hypothesis
    we have $a_{n}\rightarrow{x}_{0}$, and hence there is an $N\in\mathbb{N}$
    such that $n\in\mathbb{N}$ and $n>N$ implies
    $|a_{n}-x_{0}|<\delta$. But if $|a_{n}-x_{0}|<\delta$, then
    $|f(a_{n})-f(x_{0})|<\varepsilon$. Hence for all $n\in\mathbb{N}$ with
    $n>N$ we have $|f(a_{n})-f(x_{0})|<\varepsilon$, and therefore
    $f(a_{n})\rightarrow{f}(x_{0})$.
    \par\hfill\par
    Next, if we know that $f:[a,\,b]\rightarrow\mathbb{R}$ is continuous
    at each point, then given any convergent sequence
    $c:\mathbb{N}\rightarrow\mathbb{R}$ with $c_{n}\rightarrow{r}$,
    $r\in[a,\,b]$, we have $f(c_{n})\rightarrow{f}(r)$ by the previous
    argument. But if $x:\mathbb{N}\rightarrow[a,\,b]$ is a sequence, then by
    the Bolzano-Weierstrass theorem is a convergent subsequence $x_{k}$.
    That is, there is a strictly increasing sequence
    $k:\mathbb{N}\rightarrow\mathbb{N}$ and a real number $r\in[a,\,b]$
    such that $x_{k_{n}}\rightarrow{r}$. But by hypothesis $f$ is continuous
    at $r$, and hence $f(x_{k_{n}})\rightarrow{f}(r)$. That is,
    $f(x_{k_{n}})$ forms a convergent subsequence of $f(x_{n})$.
    \color{black}
    \newpage
    \textbf{Cauchy Sequences}
    \begin{itemize}
        \item
            Prove that if $a:\mathbb{N}\rightarrow\mathbb{R}$ and
            $b:\mathbb{N}\rightarrow\mathbb{R}$ are Cauchy sequences, and if
            there is some $r>0$ with $b_{n}>r$ for all $n\in\mathbb{N}$,
            then the quotient $c_{n}=a_{n}/b_{n}$ is also a Cauchy sequence.
        \item
            Prove that if $\alpha>0$ and $a_{n}=(n+\alpha)^{2}-n^{2}$, then
            $a:\mathbb{N}\rightarrow\mathbb{R}$ does \textbf{not} form a
            Cauchy sequence.
    \end{itemize}
    \color{blue}
    Since $b_{n}>r$ for all $n\in\mathbb{N}$, we have
    $1/b_{n}<1/r$ for all $n\in\mathbb{N}$ as well. Furthermore, since
    $a:\mathbb{N}\rightarrow\mathbb{R}$ is a Cauchy sequence, it must be
    bounded. Similarly for $b$. Let $M>0$ be a bound for both $a$ and $b$.
    Given $\varepsilon>0$, since $a$ and $b$ are Cauchy sequences, there is
    an $N\in\mathbb{N}$ such that for all
    $m,n\in\mathbb{N}$ with $m,n>N$, we have
    $|a_{n}-a_{m}|<\frac{\varepsilon}{2M/r^{2}}$ and
    $|b_{n}-b_{m}|<\frac{\varepsilon}{2M/r^{2}}$. But then:
    \begin{equation}
        \begin{array}{rcl}
            \displaystyle
            \left|
                \frac{a_{n}}{b_{n}}
                -
                \frac{a_{m}}{b_{m}}
            \right|
            \!\!
            &=&
            \!\!
            \displaystyle
            \left|
                \frac{a_{n}b_{m}-a_{m}b_{n}}{b_{n}b_{m}}
            \right|\\[1.5em]
            \!\!
            &=&
            \!\!
            \displaystyle
            \frac{1}{|b_{n}b_{m}|}
            \left|
                a_{n}b_{m}
                -
                a_{m}b_{n}
            \right|\\[1.5em]
            \!\!
            &<&
            \!\!
            \displaystyle
            \frac{1}{r^{2}}
            \left|
                a_{n}b_{m}
                -
                a_{m}b_{n}
            \right|\\[1.5em]
            \!\!
            &=&
            \!\!
            \displaystyle
            \frac{1}{r^{2}}
            \left|
                a_{n}b_{m}-a_{n}b_{n}+a_{n}b_{n}-a_{m}b_{n}
            \right|\\[1.5em]
            \!\!
            &\leq&
            \!\!
            \displaystyle
            \frac{1}{r^{2}}\left(
                \left|
                    a_{n}b_{m}-a_{n}b_{n}
                \right|
                +
                \left|
                    a_{n}b_{n}-a_{m}b_{n}
                \right|
            \right)\\[1.5em]
            \!\!
            &=&
            \!\!
            \displaystyle
            \frac{1}{r^{2}}\left(
                \left|
                    a_{n}
                \right|\,
                \left|
                    b_{m}-b_{n}
                \right|
                +
                \left|
                    b_{n}
                \right|\,
                \left|
                    a_{n}-a_{m}
                \right|
            \right)\\[1.5em]
            \!\!
            &\leq&
            \!\!
            \displaystyle
            \frac{M}{r^{2}}\left(
                \left|
                    b_{m}-b_{n}
                \right|
                +
                \left|
                    a_{n}-a_{m}
                \right|
            \right)\\[1.5em]
            \!\!
            &<&
            \!\!
            \displaystyle
            \frac{M}{r^{2}}\left(
                \frac{\varepsilon}{2M/r^{2}}
                +
                \frac{\varepsilon}{2M/r^{2}}
            \right)\\[1.5em]
            \!\!
            &=&
            \!\!
            \displaystyle
            \varepsilon
        \end{array}
    \end{equation}
    For the sequence $a_{n}=(n+\alpha)^{2}-n^{2}$, we simplify it and get:
    \begin{equation}
        a_{n}
        =
        2\alpha{n}+\alpha^{2}
    \end{equation}
    Since $\alpha>0$, $a_{n}$ becomes arbitrarily large. But Cauchy sequences
    are bounded, and hence $a$ is not a Cauchy sequence.
    \color{black}
    \newpage
    \textbf{Compactness, Open and Closed Sets}
    \begin{itemize}
        \item
            Prove that if $\mathcal{O}$ is a collection of open subsets of
            $\mathbb{R}$ (possibly empty, possibly infinite), then
            $\bigcup\mathcal{O}$ is an open subset of $\mathbb{R}$ as well.
        \item
            Prove that if $A\subseteq\mathbb{R}$ is compact, and if
            $\mathcal{C}_{n}$ are closed, non-empty, nested subsets of
            $A$ (that is, $\mathcal{C}_{n}$ is non-empty and
            $\mathcal{C}_{n+1}\subseteq\mathcal{C}_{n}$ for each
            $n\in\mathbb{N}$), then
            $\bigcap_{n=0}^{\infty}\mathcal{C}_{n}$ is non-empty.
    \end{itemize}
    \color{blue}
    If $\bigcup\mathcal{O}$ is not open, then by definition there is an
    $x\in\bigcup\mathcal{O}$ such that for all $\varepsilon>0$ there is a
    $y\in\mathbb{R}$ with $|x-y|<\varepsilon$, but $y\not\in\bigcup\mathcal{O}$.
    By definition of union, if $x\in\bigcup\mathcal{O}$, then there is some
    $\mathcal{U}\in\mathcal{O}$ such that $x\in\mathcal{U}$. But by
    hypothesis, if $\mathcal{U}\in\mathcal{O}$, then $\mathcal{U}$ is open.
    But if $x\in\mathcal{U}$ and $\mathcal{U}$ is open, then by definition
    there is some $\varepsilon>0$ such that for all $y\in\mathbb{R}$,
    $|x-y|<\varepsilon$ implies $y\in\mathcal{U}$. But then for all such
    $y$ we have $y\in\bigcup\mathcal{O}$ since
    $\mathcal{U}\subseteq\bigcup\mathcal{O}$ (by the definition of union),
    which is a contradiction. Hence $\bigcup\mathcal{O}$ is open.
    \par\hfill\par
    For the second part, for each $n\in\mathbb{N}$, define
    $\mathcal{U}_{n}=\mathbb{R}\setminus\mathcal{C}_{n}$.
    Since $\mathcal{C}_{n}$
    is closed, $\mathcal{U}_{n}$ is thus the complement of a closed set, and
    is therefore open. Suppose $\bigcap_{n=0}^{\infty}\mathcal{C}_{n}$ is
    empty. If this were true, then the collection of all $\mathcal{U}_{n}$
    form an open cover for $A$. That is, if we define $\mathcal{O}$ via
    \begin{equation}
        \mathcal{O}
        =
        \left\{\,
            \mathcal{U}_{n}\;|\;n\in\mathbb{N}\,
        \right\},
    \end{equation}
    then $A\subseteq\bigcup\mathcal{O}$. For if $x\in{A}$, then
    since we are assuming $\bigcap_{n=0}^{\infty}\mathcal{C}_{n}$ is empty,
    there must be some $n\in\mathbb{N}$ such that $x\not\in\mathcal{C}_{n}$
    (by the definition of intersection). But if $x\not\in\mathcal{C}_{n}$,
    then $x\in\mathbb{R}\setminus\mathcal{C}_{n}$, and hence
    $x\in\mathcal{U}_{n}$. That is, for all $x\in{A}$ there is an
    $n\in\mathbb{N}$ such that $x\in\mathcal{U}_{n}$, and hence
    $\mathcal{O}$ is an open cover for $A$. But $A$ is compact, and therefore
    there is a finite subcover
    $\Delta=\{\mathcal{U}_{n_{0}},\,\cdots,\,\mathcal{U}_{n_{N}}\}$.
    But the sets $\mathcal{C}_{n}$ are nested,
    $\mathcal{C}_{n+k}\subseteq\mathcal{C}_{n}$ for $k\geq{0}$, and hence
    the sets $\mathcal{U}_{n}$ are nested as well, but in reverse:
    $\mathcal{U}_{n}\subseteq\mathcal{U}_{n+k}$. Let $M$ be the largest index
    among $n_{0}$, $\dots$, $n_{N}$. Since the sets $\mathcal{U}_{n}$ are
    nested, this shows that $\{\,\mathcal{U}_{M}\,\}$ is an indeed a subcover
    for $\mathcal{O}$. That is, $A\subseteq\mathcal{U}_{M}$. But
    $\mathcal{U}_{M}=\mathbb{R}\setminus\mathcal{C}_{M}$, which means
    $A\subseteq\mathbb{R}\setminus\mathcal{C}_{M}$. But
    $\mathcal{C}_{M}\subseteq{A}$ by hypothesis. But the only way for
    $\mathcal{C}_{M}\subseteq{A}$ to be true and
    $A\subseteq\mathbb{R}\setminus\mathcal{C}_{M}$ to be true is if
    $\mathcal{C}_{M}$ is empty. But by hypothesis $\mathcal{C}_{M}$ is
    non-empty, which is a contradiction. Hence,
    $\bigcap_{n=0}^{\infty}\mathcal{C}_{n}$ is non-empty.
    \color{black}
    \newpage
    \textbf{Continuity}
    \begin{itemize}
        \item (10 Pts)
            Prove using an $\varepsilon-\delta$ argument that
            $f:[1,\,\infty)\rightarrow\mathbb{R}$ given by
            $f(x)=\frac{1}{x}$ is continuous at each $x_{0}$ in the domain.
    \end{itemize}
    \color{blue}
    For suppose $x_{0}\in[1,\,\infty)$ and $\varepsilon>0$.
    We want $|x-x_{0}|<\delta$ implies
    $|f(x)-f(x_{0})|<\varepsilon$. That is, we want:
    \begin{equation}
        \left|
            \frac{1}{x}-\frac{1}{x_{0}}
        \right|
        <
        \varepsilon
    \end{equation}
    This is equivalent to:
    \begin{equation}
        \frac{\left|x-x_{0}\right|}{\left|xx_{0}\right|}
        <
        \varepsilon
    \end{equation}
    By definition, the domain of the function is $[1,\,\infty)$, and hence
    $1\leq{x}$ and $1\leq{x_{0}}$. But then
    $1/x\leq{1}$ and $1/x_{0}\leq{1}$. Hence:
    \begin{equation}
        \frac{\left|x-x_{0}\right|}{\left|xx_{0}\right|}
        \leq
        |x-x_{0}|
    \end{equation}
    Thus, if we can make $|x-x_{0}|<\varepsilon$ true, then we automatically
    make $|f(x)-f(x_{0})|<\varepsilon$ true as well.
    Choose $\delta=\varepsilon$.
    If $x\in[1,\,\infty)$ and $|x-x_{0}|<\delta$, then:
    \begin{equation}
        \begin{array}{rcl}
            \displaystyle
            \left|f(x)-f(x_{0})\right|
            \!\!
            &=&
            \!\!
            \displaystyle
            \left|
                \frac{1}{x}-\frac{1}{x_{0}}
            \right|\\[1.5em]
            \!\!
            &=&
            \!\!
            \displaystyle
            \frac{\left|x-x_{0}\right|}{\left|xx_{0}\right|}\\[1.5em]
            \!\!
            &\leq&
            \!\!
            \displaystyle
            \left|x-x_{0}\right|\\[1.5em]
            \!\!
            &<&
            \!\!
            \displaystyle
            \delta\\[1.5em]
            \!\!
            &=&
            \!\!
            \varepsilon
        \end{array}
    \end{equation}
    Hence $|x-x_{0}|<\delta$ implies $|f(x)-f(x_{0})|<\varepsilon$.
\end{document}
